\section{Methodology}
\label{sec:method}

\subsection{Problem Formulation}
\label{sec:method:problem}

Given a biological protocol step $s$ (a text string), the task is to output a corrected version $\hat{s}^*$ and a binary flag $\hat{e} \in \{0, 1\}$ indicating whether $s$ contains an error.
The ground truth is a pair $(s^*, e)$ where $s^*$ is the expert-corrected protocol and $e$ is the true error indicator.
We evaluate agents on three metrics derived from $(\hat{s}^*, \hat{e})$: detection accuracy, exact-match correction accuracy, and semantic correction accuracy (defined in \secref{sec:method:eval}).

\subsection{The WetLabAgent System}
\label{sec:method:agent}

\wetlabagent is a modular system built around \claude~\cite{anthropic2026claude} as the reasoning backbone.
We evaluate four agent configurations of increasing sophistication.

\para{Condition 1: Vanilla LLM (\vanilla).}
The agent issues a single prompt to \claude containing the protocol step and requests a JSON-structured response with fields \texttt{has\_error} (boolean) and \texttt{corrected\_text} (string).
No external tools are called.
This serves as the baseline measuring the raw capability of \claude on protocol error correction without augmentation.

\para{Condition 2: Tool-Augmented LLM (\toolaug).}
Before calling the LLM, the agent executes two domain-specific validation tools and appends their outputs to the prompt context:
\begin{itemize}[leftmargin=*,itemsep=0pt,topsep=0pt]
    \item \textbf{\texttt{tool\_validate\_parameters}}: Uses regular expressions to extract numerical values and compares them against known biological ranges (CO$_2$: 4--6\%, incubator temperature: 35--39$^\circ$C, ROCK inhibitor dilution: 1:500--1:2000). Returns flagged anomalies.
    \item \textbf{\texttt{tool\_check\_protocol\_structure}}: Checks structural consistency between the current step and adjacent context steps (purpose, prior step, next step). Returns inconsistencies.
\end{itemize}
This tests whether domain-specific validation outputs improve LLM correction accuracy.

\para{Condition 3: Feedback Loop, 1 Round (\feedbackone).}
After \toolaug generates an initial correction $\hat{s}^*_0$, a simulated execution module applies heuristic checks to $\hat{s}^*_0$ and produces a structured observation: a pass/fail outcome and, on failure, an anomaly type (\eg \texttt{parameter\_out\_of\_range}, \texttt{structural\_inconsistency}).
This observation is appended to the conversation and the LLM generates a refined correction $\hat{s}^*_1$.

\para{Condition 4: Feedback Loop, 3 Rounds (\feedbackthree).}
The same simulated execution + LLM refinement loop runs for up to three rounds.
An early-exit condition terminates the loop when the simulated executor reports a passing outcome, yielding an average of 2.4 API calls per sample.
This models the closed-loop adaptation cycle that would occur in a physical wet-lab: generate $\to$ execute $\to$ observe $\to$ refine $\to$ repeat.

\subsection{Dataset}
\label{sec:method:data}

We use the error correction task from \bioprobench~\cite{bioprobench2025}, a benchmark of 27,000 biological protocols with 556,000 structured instances.
We evaluate on a local sample set of $n=10$ instances drawn from the full dataset.
Each instance provides: a potentially corrupted protocol step (\texttt{corrupted\_text}), the expert-corrected version (\texttt{corrected\_text}), an error type label, and a correctness flag (\texttt{is\_correct}).

The 10 samples include 6 instances with deliberate errors and 4 correct instances.
Error types are: parameter (2 instances, \eg CO$_2$ concentration), reagent (2 instances, \eg dilution factor), and operation (2 instances, \eg centrifuge speed).
For correct instances, the LLM input is \texttt{corrected\_text} and the expected output is \texttt{has\_error=False}.
\Tabref{tab:dataset} summarizes the sample composition.

\begin{table}[t]
\centering
\caption{Sample composition of our \bioprobench evaluation set ($n=10$). Error types cover parameter, reagent, and operation categories; four samples contain no error.}
\label{tab:dataset}
\small
\begin{tabular}{@{}lcc@{}}
\toprule
\textbf{Error Type} & \textbf{Count} & \textbf{Example} \\
\midrule
Parameter & 2 & Wrong CO$_2$ concentration (10\% instead of 5\%) \\
Reagent   & 2 & Wrong ROCK inhibitor dilution (1:100 instead of 1:1000) \\
Operation & 2 & Wrong centrifuge speed (1600$\times$g instead of 800$\times$g) \\
Correct   & 4 & No error (LLM should output \texttt{has\_error=False}) \\
\midrule
\textbf{Total} & 10 & \\
\bottomrule
\end{tabular}
\end{table}

\subsection{Evaluation Metrics}
\label{sec:method:eval}

We use three primary metrics:

\para{Error Detection Accuracy.}
Binary accuracy: whether the agent correctly classifies the protocol as containing an error or not.
For the 6 corrupted samples, detection is a true positive iff \texttt{has\_error=True}.
For the 4 correct samples, detection is a true negative iff \texttt{has\_error=False}.

\para{Exact Correction Accuracy.}
Binary accuracy: whether the agent's corrected text exactly matches the ground truth \texttt{corrected\_text} after whitespace normalization.
This is the strictest evaluation criterion, penalizing any deviation from the reference text including formatting changes.

\para{Semantic Correction Accuracy.}
Binary accuracy: whether the agent's output contains the correct primary fix, even if it also modifies other parts of the text.
Evaluated using heuristic string matching for the key changed value (\eg presence of ``5\%'' in a CO$_2$ correction task).
This better captures the scientific correctness of the correction.

\para{Secondary metrics.}
We also report average API latency (seconds per sample) and average API call count per sample to quantify the efficiency cost of each condition.

\subsection{Implementation Details}
\label{sec:method:impl}

All experiments use \claude with temperature 0.0 for determinism, maximum 512 output tokens, and random seed 42 for all stochastic operations.
The Anthropic Python SDK (v0.84.0) is used for API calls.
Experiments ran on a CPU-only Linux system (kernel 6.12.72), with a total of approximately 50 API calls across all conditions.
Statistical comparisons use McNemar's test for binary outcomes with Cohen's $h$ as the effect size measure.
